% Generated draft: requires mathematical and source review.
\documentclass[11pt]{article}
\usepackage[margin=25mm]{geometry}
\usepackage{fontspec}
\setmainfont{Latin Modern Roman}
\usepackage{amsmath,amssymb,mathtools,unicode-math}
\setmathfont{Latin Modern Math}
\usepackage{enumitem,needspace,xurl,hyperref}
\setlength{\parindent}{0pt}
\setlength{\parskip}{5pt}
\setlength{\emergencystretch}{3em}
\newcommand{\N}{\mathbb{N}}
\newcommand{\Z}{\mathbb{Z}}
\newcommand{\Q}{\mathbb{Q}}
\newcommand{\R}{\mathbb{R}}
\newcommand{\C}{\mathbb{C}}
\newcommand{\F}{\mathbb{F}}
\newcommand{\A}{\mathbb{A}}
\newcommand{\PP}{\mathbb{P}}
\newcommand{\E}{\mathbb{E}}
\newcommand{\eps}{\varepsilon}
\newcommand{\dd}{\,\mathrm{d}}
\begin{document}
\section*{UnsolvedMath ID: 30006509}
\textbf{Source problem code:} catalogue entry
\subsection{Definitions and mathematical statement}
Let \(D=\{z\in\C:|z|<1\}\). The supplied excerpt considers a Brownian loop soup \(\mathcal L_c\) in \(D\), a Poisson collection of Brownian loops with intensity parameter \(c>0\), in a convention whose critical intensity is \(c=1\). It does not specify the normalization of Brownian loop measure.

For a conditional formulation, connect two loops when their traces intersect, and define a cluster as an equivalence class under finite chains of such connections. For a cluster \(\mathcal C\), put
\[
K_{\mathcal C}=\overline{\bigcup_{\ell\in\mathcal C}\operatorname{tr}(\ell)},\qquad
\partial_{\mathrm{out}}K_{\mathcal C}=\partial U_\infty(K_{\mathcal C}),
\]
where \(U_\infty(K_{\mathcal C})\) is the unbounded connected component of \(\C\setminus K_{\mathcal C}\). This states an explicit closure and outer-boundary convention; the excerpt does not itself supply these definitions.

The term ``double point'' is not defined in the supplied excerpt. Two distinct candidate meanings can be made precise. Represent each loop \(\ell\) with duration \(T_\ell>0\), taking times in \([0,T_\ell)\) so that periodic endpoint identification alone does not count as a repeated visit. Define
\[
\begin{aligned}
\mathcal D_{\mathrm{same}}(\mathcal C)
&=\partial_{\mathrm{out}}K_{\mathcal C}\cap
\bigcup_{\ell\in\mathcal C}
\{z:\exists\,0\le s<t<T_\ell,\ \ell(s)=\ell(t)=z\},\\
\mathcal D_{\mathrm{distinct}}(\mathcal C)
&=\partial_{\mathrm{out}}K_{\mathcal C}\cap
\bigcup_{\substack{\ell,\ell'\in\mathcal C\\\ell\ne\ell'}}
\bigl(\operatorname{tr}(\ell)\cap\operatorname{tr}(\ell')\bigr).
\end{aligned}
\]
These respectively describe repeated visits by one loop and visits by two distinct loops. Their union is another possible interpretation; the record does not establish which interpretation is intended, or whether multiplicity must be exactly two.

The excerpt motivates an existence question for \(c\in(0,1]\), while the supplied literature assessment singles out \(c=1\). For any specified choice \(\mathcal D\), one precise critical existence question would be whether
\[
\mathbb P\!\left(\exists\text{ a cluster }\mathcal C\text{ of }\mathcal L_1:
\mathcal D(\mathcal C)\ne\varnothing\right)>0.
\]
Its negation is almost-sure emptiness for every cluster. The excerpt does not specify whether its intended existence assertion is this positive-probability assertion, almost-sure existence, or an assertion about a selected cluster or only outermost clusters. Consequently, this display is a conditional formulation, not a recovered definitive statement.

There is also a substantive conflict in the supplied record. The dated literature assessment labels critical existence unresolved, whereas the supplied statement-verification note says that the official report subsequently proves nonexistence of the double points it discusses and that this extraction is therefore not an open problem. The extracted text ends at ``We are particularly interested'' and omits that continuation. No unique open problem or fully specified nonexistence theorem can therefore be asserted from the supplied evidence.
\subsection{Short English statement}
The record concerns points visited more than once on outer boundaries of Brownian loop-soup clusters, especially at critical intensity \(c=1\). The excerpt is truncated, leaves the meaning of a double point and the probability quantifiers unspecified, and has conflicting supplied status assessments. The formulas give conditional interpretations requiring human mathematical review.
\subsection{Sources}
\begin{itemize}
\item Local registry: \texttt{data/unsolvedmath/problems.json}, permanent UnsolvedMath ID \(30006509\), exact source code \texttt{catalogue entry}; supplied title: ``Reported Nonexistence of Critical Loop-Soup Double Points.'' The record belongs to the Oberwolfach Reports Open Problems collection and gives proposed year \(2025\), with no named proposer.
\item Supplied report citation: \emph{Functional Inequalities: Geometric Calculus meets Stochastic Analysis} (2026), Oberwolfach Reports, pp.~2917--2994, DOI \texttt{10.4171/owr/2025/54}. \url{https://doi.org/10.4171/owr/2025/54}. The citation year and DOI year are retained as supplied. The report was not independently consulted.
\item Supplied literature reference: Gao--Li--Qian, \emph{Multiple points on Brownian loop-soup cluster boundaries}. \url{https://arxiv.org/abs/2205.11468}.
\item Supplied literature reference: Annals of Probability, \emph{Multiple points on Brownian loop-soup cluster boundaries}. \url{https://doi.org/10.1214/25-AOP1765}.
\item Literature assessment checked at \texttt{2026-08-22}: the supplied curation labels the problem open and says that subcritical double points are understood while existence at critical intensity one remains unresolved. This is supplied dated curation, not an independently verified literature assessment.
\item The supplied statement-verification field instead says that the official report motivates the critical question and then proves nonexistence, so that the extraction is not an open problem. The record marks its statement status \texttt{unrecoverable} and supplies no clean statement. This verification claim is preserved as source metadata and has not been independently checked.
\item Historical attribution in the excerpt: Lawler and Werner introduced Brownian loop soup, cited there as [14]; Kiefer and M\"orters are credited with a Brownian-frontier double-point dimension result [8]; Duplantier and Kwon with a disconnection-exponent conjecture [4]; and Lawler, Schramm and Werner with its proof using SLE [10, 11, 13, 12]. The excerpt also cites [20], [19, Lemma 4], and [1] without supplying their bibliographic entries.
\end{itemize}
\end{document}
